\section{Discussion \& Limitations}
\label{sec:discussion}

The intended benefit of monodromy is semantic stability under shared continuous
adaptation. It protects an executable sheet action over a certified region,
while leaving geometric implementation choices available for learning. Whether
this freedom improves acquisition or resource use depends on the task and on
the cost of finding, executing, and maintaining the region. An ordinary
constraint method may exploit exactly the same freedom; the comparison must
therefore measure a residual computational benefit rather than attach value
to a topological description alone.

Several interfaces remain consequential. Stable operations can be called
incorrectly after an encoder or router changes. The native relation and
finite-length soft condition address this distinction, but checking them may
be expensive or conservative. The current integrated diagnostic uses a supplied
finite covering, small state space, bounded programs, and a supplied old
capability. It does not establish autonomous substrate discovery, unbounded
capability growth, or a general-purpose learning architecture. Reversible
lifts broaden logical expressivity only by paying for auxiliary state and
cleanup; a large configuration space is not itself a large parameter count.

A useful architecture must preserve the same core across new computational
and real-data tasks, with honest costs for input adaptation and memory. The
primary benchmark therefore waits for a validated complete Block and a frozen
implementation. Existing negative results, ordinary-method correspondences,
and failed training runs remain evidence for developing that Block, rather
than being removed from its history or treated as conclusions about all
topological learning systems.
